% Modernised reading — fonds Alexandre Grothendieck, shelfmark 134-1, pages
% 1-6 (the whole folder; pages 1 and 2 are the cover and a blank leaf, so the
% mathematics is pages 3 to 6).
%
% This is an INTERPRETATION, not a transcription. It re-reads the pages in
% today's notation and vocabulary, states the hypotheses the manuscript leaves
% implicit, and drops the critical apparatus so that the mathematics reads
% straight through. Everything the brackets and underlines carried — a doubtful
% reading, a notation he used differently, a missing passage — moves into a
% footnote.
%
% It is held to being mathematically correct as it stands, not merely faithful
% to the page. Where the manuscript is loose or elliptical, the true statement
% is what appears here, and the footnote says what the page has. In any dispute
% of substance the transcription governs: whoever cites Grothendieck must cite
% batch-01.fr.tex.
%
% Scope: a SINGLE document for the whole folder, the convention every reading
% in this repository follows. The shelfmark holds one batch, so pages 1-6 are
% the folder entire and no part of it is read in another file. The folder was
% read whole — all six pages of batch-01.fr.tex — before any of this was
% written.
%
% Four departures from the page are worth naming here, because they are what a
% reader would otherwise trip on.
%
% 1. The letter never defines E'(M,N). It says only what E'(M,N) is supposed to
%    express, heuristically — the strict 2-Picard stack made of not-necessarily
%    -strict Picard stacks. A modernised reading cannot leave the principal
%    object undefined, so a definition is supplied: E'(M,N) is the truncated
%    mapping spectrum computed over the sphere rather than over Z. That
%    definition is ours, it is flagged as ours, and everything the letter
%    asserts about E' is then a theorem rather than a heuristic.
%
% 2. Page 4 says E'(M,N) has "les invariants H^i ... ceux de E(M,N) en degré
%    i+2". Read literally that is false — the right-hand side vanishes for
%    i > 0. What the triangle gives is: H^i(E') = H^i(E) for i <= 1, and in
%    degree 2 the exact sequence (S). The correct statement is what appears
%    below, with the page's own in a footnote.
%
% 3. The direction of the triangle's arrows is what makes (S) come out rather
%    than a different four-term sequence, and it is checked here rather than
%    copied: the map runs from E to E', the forgetful map from Z-linear to
%    sphere-linear, and the third term is its cofibre.
%
% 4. The obstruction group. The page says the obstruction lies in
%    Ext^3(X; M, N). What the triangle gives is the hypercohomology group
%    H^3(X, E(M,N)), which INJECTS into Ext^3(X; M, N) — a subgroup, not the
%    whole. The finer statement is given, with the page's in a footnote.
%
% Pass: Opus 5 (claude-opus-5), 2026-08-24 — first pass, unchecked against the
% transcription by a human.
\documentclass[11pt,a4paper]{article}
% Le préambule commun à toutes les transcriptions du fonds Grothendieck.
%
% Il tient en une idée : **l'incertitude se compose, elle ne se résout pas.**
% Une transcription qui lisse un mot douteux a détruit la seule chose pour
% laquelle on la faisait — un lecteur doit pouvoir savoir, à chaque endroit,
% ce qui était sur le feuillet et ce qui a été ajouté. Les six macros qui
% suivent sont donc l'appareil critique, et non de la décoration.
%
% Les mêmes commandes sont reconnues par `scripts/render.mjs`, qui produit la
% vue de lecture du site. Une macro ajoutée ici doit l'être aussi là-bas,
% faute de quoi le rendu échouera bruyamment — ce qui est voulu : un rendu
% silencieusement faux est pire qu'un rendu absent.

\ProvidesPackage{grothendieck}[2026/08/08 Transcriptions du fonds Grothendieck]

\usepackage{amsmath,amssymb,amsthm}
\usepackage{mathrsfs}
\usepackage{stmaryrd}
% Les diagrammes commutatifs sont partout dans ce fonds : c'est la langue
% même de la géométrie algébrique de Grothendieck, et une transcription qui
% les rendrait en prose ne transcrirait rien.
\usepackage{tikz-cd}
% Français par défaut — c'est la langue des feuillets — mais l'anglais est
% chargé aussi : la traduction et le résumé sont des documents anglais, et la
% typographie française (espaces avant les deux-points) y serait une faute.
% Ces documents basculent par \selectlanguage{english} en tête de corps.
\usepackage[english,french]{babel}
\usepackage{fontspec}
\usepackage{geometry}
\geometry{a4paper,margin=25mm,marginparwidth=28mm}
\usepackage{marginnote}
\usepackage{xcolor}
% ulem, not soul. `soul`'s \st and \ul are built on a character-by-character
% scan that XeLaTeX's font machinery breaks: the document fails with an
% undefined control sequence rather than degrading. ulem does the same two jobs
% and is Unicode-engine safe. `normalem` stops it hijacking \emph, which the
% transcriptions use for Grothendieck's own emphasis.
\usepackage[normalem]{ulem}
\usepackage{hyperref}
\hypersetup{colorlinks=true,linkcolor=black,urlcolor=black,pdfborder={0 0 0}}

% --- Métadonnées du lot -----------------------------------------------------
% Reprises telles quelles par le script de rendu, qui les affiche en tête de
% la vue de lecture. Elles fixent ce que le fichier prétend couvrir : sans
% elles, une transcription flotte sans point d'attache dans un fonds de seize
% mille feuillets.

\newcommand{\folder}[1]{\gdef\grothFolder{#1}}
\newcommand{\batch}[1]{\gdef\grothBatch{#1}}
\newcommand{\leaves}[2]{\gdef\grothFirst{#1}\gdef\grothLast{#2}}
\newcommand{\foldertitle}[1]{\gdef\grothFoldertitle{#1}}
\gdef\grothFolder{?}\gdef\grothBatch{?}\gdef\grothFirst{?}\gdef\grothLast{?}\gdef\grothFoldertitle{}

% --- Le résumé --------------------------------------------------------------
%
% Il ouvre la lecture modernisée, sous le titre « Résumé », et il est composé
% en petit. Ce n'est pas un ornement : c'est le seul passage du document qui
% ne soit pas des mathématiques, et un lecteur doit pouvoir le reconnaître
% avant de l'avoir lu — puis le sauter s'il connaît déjà le sujet, ou s'y
% arrêter s'il ne le connaît pas. Le filet qui le suit marque la reprise du
% corps.

\newenvironment{resume}
  {\small\setlength{\parskip}{0.45em}}
  {\par\medskip\hrule\medskip}

% --- Mots-clés ---------------------------------------------------------------
%
% En anglais, en clôture du résumé de la lecture modernisée : le vocabulaire
% moderne sous lequel chercher ce que les feuillets construisent. Le manifeste
% du site (`npm run manifest`) les extrait pour étiqueter la cote entière —
% cette ligne est la source unique des tags, il n'y a pas de fichier à côté.
\newcommand{\keywords}[1]{\par\smallskip\noindent
  {\footnotesize\textsc{Keywords} --- \textit{#1}}\par}

% --- L'appareil critique ----------------------------------------------------

% Le numéro de feuillet, en marge. C'est l'ancre qui permet de régler un
% désaccord sur une formule en regardant *un* feuillet plutôt qu'une chemise
% de six cents. Le site s'en sert aussi pour tourner les pages du fac-similé
% quand on fait défiler la transcription.
\newcommand{\leaf}[1]{%
  \marginnote{\footnotesize\color{gray!70}\textsf{#1}}%
  \label{leaf:#1}\ignorespaces}

% Illisible : on ne devine pas. Un mot inventé qui se lit comme les autres est
% le pire résultat possible.
\newcommand{\ill}{\textcolor{red!60!black}{[\,\ldots\,]}}

% Lecture douteuse : le mot est donné, et signalé comme incertain.
\newcommand{\uncertain}[1]{\uline{#1}}

% Ajout de l'éditeur — une lettre manquante, un mot sous-entendu.
\newcommand{\add}[1]{\textcolor{blue!50!black}{[#1]}}

% Biffé par Grothendieck lui-même. Ce qu'il a rayé fait partie du document :
% une rature dit souvent où la pensée a changé de direction.
\newcommand{\struck}[1]{\sout{#1}}

% Note du transcripteur, sur l'état matériel du feuillet.
\newcommand{\note}[1]{\par\smallskip{\small\color{gray!60!black}\textit{[#1]}}\par\smallskip}

% Note marginale de Grothendieck, distincte du corps du texte.
\newcommand{\marginal}[1]{\par\smallskip\noindent\hspace{1em}%
  {\small\color{gray!60!black}#1}\par\smallskip}

% --- Titre ------------------------------------------------------------------

\AtBeginDocument{%
  \begin{center}
    {\large\bfseries Fonds Alexandre Grothendieck --- cote n\textsuperscript{o}~\grothFolder}\\[2pt]
    {\itshape\grothFoldertitle}\\[4pt]
    {\small feuillets \grothFirst--\grothLast\ (lot \grothBatch)}\\[2pt]
    {\footnotesize\color{gray!60!black}Transcription établie d'après le fac-similé numérisé
     par l'Université de Montpellier.\\
     Les lectures incertaines sont soulignées, les illisibles notées [\,\ldots\,],
     les ajouts entre crochets.}
  \end{center}
  \vspace{1em}}

\endinput


\folder{134-1}
\pages{1}{6}
\dating{1974-1984}
\watermark{Édition de démonstration\\interprétation personnelle de l'œuvre}
\foldertitle{Lettres Breen 75,76 [correspondance avec Larry Breen, Pierre Deligne et Daniel Quillen] : lettres et copies de lettre (1974-1976, 1984) — lecture modernisée du dossier entier}

\begin{document}

\section*{Résumé}

\begin{resume}
Ce dossier contient une seule lettre, de quatre pages, écrite de Lodève et
adressée à Pierre Deligne. Grothendieck y demande si une petite chose qu'il
vient de remarquer est connue. La chose est celle-ci : lorsqu'on remplace,
dans un groupe commutatif, l'égalité par un isomorphisme qu'il faut choisir,
il reste exactement un invariant de plus qu'on ne l'attendait, et cet
invariant est un signe.

Voici d'où il vient. Dans un groupe commutatif, $xy$ et $yx$ sont le même
élément et il n'y a rien à dire. Supposons maintenant que les « éléments »
soient les objets d'une catégorie et que l'égalité ait été remplacée par
l'isomorphisme : il faut alors se donner, pour chaque couple d'objets, un
isomorphisme précis $c_{x,y}$ de $x \otimes y$ sur $y \otimes x$. Prenons-le
avec $y = x$. On obtient un automorphisme de $x \otimes x$, et rien n'oblige
cet automorphisme à être l'identité ; le seul axiome disponible, qui dit
qu'échanger deux fois ne fait rien, oblige seulement son carré à valoir
l'identité. Cet automorphisme est donc d'ordre au plus $2$, et il n'est pas
toujours trivial : en algèbre graduée, échanger une droite de degré impair
avec elle-même donne $-1$, et c'est toute la règle des signes.

On veut classer ces objets. Deux invariants sont évidents : le groupe $M$ des
objets à isomorphisme près, et le groupe $N$ des automorphismes de l'objet
unité. La question est de savoir ce qu'il y a d'autre, et elle a deux
réponses selon qu'on interdit ou qu'on autorise le signe. Si on l'interdit
--- on dit alors que l'objet est \emph{strict} --- la réponse était connue :
c'est un groupe $\mathrm{Ext}^2(M,N)$, l'analogue en un cran plus haut de la
classification des extensions de groupes par une classe de cohomologie. Si on
l'autorise, il faut ajouter le signe, c'est-à-dire un homomorphisme de $M$
dans les éléments d'ordre $2$ de $N$.

L'idée de la lettre est de mettre les deux réponses dans un même objet. À
$M$ et $N$ elle associe deux complexes : l'un, $E(M,N)$, dont la cohomologie
en degré $2$ classe les objets stricts ; l'autre, $E'(M,N)$, dont la
cohomologie en degré $2$ classe tous les objets. Elle affirme qu'ils sont
reliés par un triangle distingué dont le troisième terme est le signe tout
seul, placé exactement en degré $2$. On en tire une suite exacte courte $(S)$
qui dit : classer un objet quelconque, c'est classer un objet strict et
choisir un signe --- ni plus, ni moins. C'est la même forme qu'une suite de
coefficients universels, et elle joue le même rôle : elle mesure d'un seul
coup ce qu'on perd en travaillant dans le monde plus rigide.

Grothendieck dit ensuite savoir démontrer que tout signe se réalise
effectivement, ce qui n'a rien d'évident : l'obstruction vit a priori dans un groupe de cohomologie de degré $3$, et un argument « universel » la tue. Il termine par
un exemple qui est la raison d'être de tout le reste : lorsque $M$ et $N$ sont
les groupes $K^{0}$ et $K^{1}$ d'un anneau, le signe est l'échange de deux
copies d'un même module, c'est-à-dire une matrice de permutation de
déterminant $-1$ --- de sorte que la K-théorie fournit, gratuitement, une
section canonique de la suite $(S)$.

Où cela mène-t-il ? Une catégorie de Picard est aujourd'hui la même chose
qu'un spectre n'ayant que deux groupes d'homotopie, en degrés $0$ et $1$ ; le
signe est la multiplication par l'élément $\eta$ de la sphère, et la classe
$\mathrm{Ext}^2$ est le premier invariant de Postnikov. Le triangle de la
lettre est alors la comparaison entre « linéaire sur les entiers » et
« linéaire sur la sphère », dont le premier écart est exactement d'ordre $2$.
La question finale de la lettre --- y a-t-il des variantes supérieures ? ---
reçoit dans ce langage une réponse de principe : oui, une par étage de la
tour de Postnikov de la sphère. Les noms à chercher sont \emph{champ de
Picard}, \emph{type d'homotopie stable}, \emph{invariant de Postnikov},
\emph{carré de Steenrod $Sq^{2}$}.

Reste la manière. La chose n'arrive pas au terme d'un programme : elle arrive
parce qu'il prépare, faute de mieux, un petit séminaire d'algèbre sur la thèse
de Hoàng Xuân Sính, et qu'en la transposant dans le langage des champs il
« tombe sur le truc suivant ». La lettre est entièrement faite d'hésitations
calibrées --- « qui pour l'instant reste heuristique », « sauf erreur », « ce
qui pour moi ne fait guère de doute » --- et le mot « splittée », écrit puis
barré, est remplacé par un énoncé plus faible et exact. Elle ne conclut pas :
elle demande. Et dans la marge de la première page, d'une autre main, Deligne
la fait suivre à Larry Breen avec ce mot : « j'imagine que tu sais ce que le
maître demande, mais n'ai pu néanmoins lui répondre ».

\keywords{Picard stack, stable homotopy 1-type, Postnikov invariant, eta multiplication, sphere spectrum, algebraic K-theory}
\end{resume}

\pagerange{3}{3}
\section*{La lettre et son occasion (page 3)}

La lettre est datée « Lodève le 27 », le mois et l'année ayant été
surchargés puis perdus à la photocopie.\footnote{Le dossier est intitulé
« Lettres Breen 75, 76 » et l'inventaire date l'ensemble de 1974-1984. La
thèse de Hoàng Xuân Sính, dont la lettre parle au futur comme d'un sujet de
séminaire, a été soutenue en 1975. Ces trois indications sont compatibles avec
1975 ou 1976 ; aucune ne date la lettre, et cette lecture ne la date pas. Les
feuillets portent par ailleurs 60 à 64 au crayon de l'archiviste --- c'est la
numérotation de la cote 134 prise entière --- tandis que les numéros employés
ici, 1 à 6, sont ceux de la partie 134-1.}

Son occasion est un contretemps d'enseignement : n'ayant peut-être pas de
cours de premier cycle à assurer, Grothendieck envisage de faire à la place un
petit séminaire d'algèbre sur « les fourbis de Mme Sinh », c'est-à-dire sur la
thèse de Hoàng Xuân Sính, éventuellement transposés dans le langage des
champs.\footnote{Le tapuscrit de cette thèse et les notes qui l'accompagnent
forment le dossier 135 du même fonds. La phrase où il explique pourquoi il
est libre de son trimestre est en partie illisible sur la page ; seul le
projet de séminaire est assuré.} C'est en la transposant qu'il rencontre ce
que la lettre expose, et qu'il présente d'emblée comme heuristique.

La marge gauche de cette même page porte, à la verticale et d'une autre main,
un mot de Pierre Deligne à Larry Breen : il fait suivre la lettre, dit
imaginer que Breen sait ce que « le maître » demande, et avoue n'avoir pu lui
répondre. Le dossier est donc l'exemplaire de Breen.

\pagerange{3}{5}
\section*{Champs de Picard épinglés par deux faisceaux (pages 3 à 5)}

Fixons un topos $X$. Un \emph{champ de Picard} sur $X$ est un champ en
groupoïdes $P$ muni d'une loi $\otimes : P \times P \to P$, d'une contrainte
d'associativité et d'une contrainte de commutativité
$c_{x,y} : x \otimes y \xrightarrow{\ \sim\ } y \otimes x$ vérifiant les
axiomes de cohérence, tel que $P$ soit localement non vide et que
$x \mapsto x \otimes y$ soit une équivalence pour tout $y$.\footnote{C'est la
définition de Deligne, SGA 4, Exposé XVIII, 1.4. La lettre l'emploie sans la
rappeler, ce qui est normal entre les deux correspondants ; elle est rappelée
ici parce que tout ce qui suit en dépend, et notamment la clause de symétrie
$c_{y,x} \circ c_{x,y} = \mathrm{id}$, qui est ce qui fait de l'invariant
$\sigma$ un élément d'ordre $2$.}

Un tel champ a deux invariants immédiats :
\[ \pi_0(P) = \text{faisceau des classes d'objets}, \qquad
   \pi_1(P) = \underline{\mathrm{Aut}}(0). \]
Pour tout objet $y$, la translation par $y$ identifie canoniquement
$\underline{\mathrm{Aut}}(y)$ à $\pi_1(P)$, de sorte que $\pi_1(P)$ mesure les
automorphismes de n'importe quel objet.

Soient maintenant $M$ et $N$ deux faisceaux abéliens sur $X$. Dire qu'un champ
de Picard $P$ est \emph{épinglé} par $M$ et $N$, c'est le donner avec deux
isomorphismes $\pi_0(P) \simeq M$ et $\pi_1(P) \simeq N$ ; les équivalences
entre tels objets sont celles qui respectent les épinglages.\footnote{Le mot
« épinglé » est celui de la lettre ; il n'y est pas défini. La définition
ci-dessus est la seule sous laquelle les énoncés de la lettre sont vrais :
sans les isomorphismes, les classes d'équivalence ne formeraient pas un
$\mathrm{Ext}^2$ mais un quotient de celui-ci par l'action des automorphismes
de $M$ et de $N$.}

\subsection*{L'invariant de symétrie}

Soit $P$ un champ de Picard et $L$ une section de $P$ au-dessus d'un objet $U$
de $X$. La contrainte de commutativité prise en $y = x = L$ donne un
automorphisme
\[ c_{L,L} \in \underline{\mathrm{Aut}}(L \otimes L) \simeq \pi_1(P) . \]
L'axiome de symétrie donne $c_{L,L}^{\,2} = \mathrm{id}$, et l'axiome
d'associativité (l'hexagone) donne
$c_{L \otimes L', L \otimes L'} = c_{L,L} \cdot c_{L',L'}$, puisque les deux
termes croisés $c_{L,L'}$ et $c_{L',L}$ sont inverses l'un de l'autre. Cet
invariant ne dépend donc que de la classe de $L$ dans $\pi_0(P)$, il est
additif, et il est tué par $2$ : c'est un homomorphisme
\[ \sigma(P) \,:\, M \longrightarrow {}_2 N , \]
où ${}_2 N$ désigne le sous-faisceau des sections de $N$ tuées par
$2$.\footnote{Page 5 : « $\sigma$ étant l'invariant associé à toute section
$L$ d'un champ de Picard, la symétrie de $L \otimes L$, interprétée comme
section de ${}_2N$ ». Le calcul qui montre que c'est un homomorphisme n'est
pas sur la page ; il est classique et donné ici pour que l'énoncé de $(S)$ ait
un sens. On notera l'identification $\underline{\mathrm{Hom}}(M, {}_2N) =
\underline{\mathrm{Hom}}(M/2M, N)$, dont on se servira plus bas : un
homomorphisme à valeurs dans ${}_2N$ est exactement un homomorphisme qui se
factorise par $M/2M$.}

Un champ de Picard est dit \emph{strict} lorsque $c_{x,x} = \mathrm{id}$ pour
tout $x$. Vu ce qui précède, cela équivaut à $\sigma(P) = 0$. Le signe est
donc \emph{le} défaut de stricte commutativité, et c'est le seul.

\subsection*{Les deux classifications}

Le cas strict est le cas connu. Le théorème de Deligne identifie la
$2$-catégorie des champs de Picard stricts sur $X$ à la catégorie dérivée
$D^{[-1,0]}(X)$, un champ $P$ correspondant à un complexe $K$ avec
$H^{-1}(K) = \pi_1(P)$ et $H^{0}(K) = \pi_0(P)$.\footnote{SGA 4, Exposé
XVIII, 1.4.15. Le séminaire a paru en 1972-1973, donc avant la lettre, et cet
exposé est de Deligne lui-même : c'est le destinataire qui avait démontré le
résultat, ce qui explique que la lettre le tienne pour acquis sans le nommer.} Un tel
complexe épinglé par $M$ et $N$ est un dévissage
$N[1] \to K \to M \xrightarrow{\,k\,} N[2]$, entièrement déterminé par sa
classe $k$, d'où la classification par
\[ \underline{\mathrm{Ext}}^2(M,N) . \]

Le cas général est celui que la lettre veut atteindre. Grothendieck introduit
d'abord le complexe qui porte le cas strict,
\[ E(M,N) \;=\; \tau_{\leq 2} \, R\underline{\mathrm{Hom}}(M,N) , \qquad
   \underline{H}^i\bigl(E(M,N)\bigr) =
   \begin{cases}
     \underline{\mathrm{Ext}}^i(M,N) & 0 \leq i \leq 2, \\
     0 & \text{sinon,}
   \end{cases} \]
puis pose l'existence d'un second complexe $E'(M,N)$ dont le
$\underline{H}^2$, noté $P(M,N)$, est le faisceau des classes d'équivalence de
champs de Picard \emph{non nécessairement stricts} épinglés par $M$ et
$N$.\footnote{Plus précisément : $P(M,N)$ est le faisceau associé au
préfaisceau $U \mapsto \{\text{champs de Picard sur } X/U \text{ épinglés par
} M, N\}/\!\simeq$. La lettre dit « le faisceau des données à équivalence près
des champs de Picard épinglés par $M$, $N$ » ; la faisceautisation est
implicite et elle est nécessaire, la formation des classes d'équivalence
n'étant pas de nature locale.}

La lettre ne définit pas $E'(M,N)$ : elle dit seulement ce qu'il doit
exprimer, à savoir le $2$-champ de Picard strict formé des $1$-champs de
Picard non nécessairement stricts, « en admettant que la théorie pour les
$1$-champs de Picard stricts s'étend aux $2$-champs de Picard stricts (ce qui
pour moi ne fait guère de doute) ». Cette lecture en donne une définition,
qui est la nôtre, et sous laquelle tout ce que la lettre affirme devient
démontrable :
\[ E'(M,N) \;=\; \tau_{\leq 2} \, R\underline{\mathrm{Hom}}_{\mathbb{S}}(HM, HN), \]
le complexe de morphismes calculé dans les faisceaux de spectres sur $X$,
c'est-à-dire au-dessus de la sphère $\mathbb{S}$ et non plus au-dessus de
$\mathbb{Z}$.\footnote{Cette définition est la nôtre, et elle est anachronique
au sens où le langage n'était pas disponible sous cette forme. Ce qu'elle
formalise ne l'est pas : un champ de Picard sur $X$ est la même chose qu'un
faisceau de spectres connectifs tronqué en degré $1$, les champs stricts
étant ceux qui sont des modules sur $H\mathbb{Z}$. La justification qu'elle
donne à $(T)$ est exposée à la section suivante. Le lecteur qui préfère s'en
tenir à la lettre peut lire $E'(M,N)$ comme un objet caractérisé par
$\underline{H}^2(E') = P(M,N)$ et par le triangle $(T)$ lui-même, ce qui est
exactement le statut heuristique que la lettre lui donne.}

\pagerange{3}{4}
\section*{Le triangle $(T)$ (pages 3 et 4)}

Le morphisme d'oubli, qui à un champ de Picard strict associe le champ de
Picard sous-jacent --- ou, dans le langage ci-dessus, qui à une application
$H\mathbb{Z}$-linéaire associe l'application $\mathbb{S}$-linéaire
sous-jacente --- fournit une flèche $E(M,N) \to E'(M,N)$. La lettre affirme
qu'elle s'insère dans un triangle distingué canonique
\begin{tikzcd}
& \underline{\mathrm{Hom}}(M, {}_2 N)[-2] \arrow[dl, "{[1]}"] & \\
E(M, N) \arrow[rr] & & E'(M, N) \arrow[ul]
\end{tikzcd}
\[ (T) \qquad E(M,N) \longrightarrow E'(M,N) \longrightarrow
   \underline{\mathrm{Hom}}(M, {}_2 N)[-2] \xrightarrow{\ [1]\ } \]

Le sens des flèches n'est pas indifférent, et il est ici celui de la page : la
flèche part de $E$ et non de $E'$.\footnote{C'est ce sens, et lui seul, qui
produit la suite $(S)$. Pris à l'envers, le même triangle donnerait la suite à
quatre termes $\underline{\mathrm{Ext}}^1 \to \underline{\mathrm{Hom}}(M,{}_2N)
\to \underline{H}^2(E') \to \underline{\mathrm{Ext}}^2 \to 0$, qui n'est pas
ce que la page écrit et qui ne dit pas ce que la page veut dire. La page
dessine le triangle sans indiquer le décalage sur la flèche de retour ; le
« $[1]$ » porté ci-dessus est le nôtre.}

\subsection*{Pourquoi il est vrai}

Le calcul, dans le langage de la définition adoptée, tient en trois lignes.
Comme $HN$ est un module sur $H\mathbb{Z}$, l'adjonction donne
\[ R\underline{\mathrm{Hom}}_{\mathbb{S}}(HM, HN) \;=\;
   R\underline{\mathrm{Hom}}_{H\mathbb{Z}}(H\mathbb{Z} \wedge HM, HN), \]
et la flèche $E \to E'$ est la restriction le long du morphisme d'action
$\mu : H\mathbb{Z} \wedge HM \to HM$, qui est bien $H\mathbb{Z}$-linéaire. Il
suffit donc de connaître la fibre de $\mu$. Les groupes d'homotopie de
$H\mathbb{Z} \wedge HM$ se calculent à partir de l'homologie entière du
spectre d'Eilenberg--MacLane, $H_*(H\mathbb{Z};\mathbb{Z}) = \mathbb{Z}, 0,
\mathbb{Z}/2, \ldots$ en degrés $0, 1, 2$ : on trouve
\[ \pi_0 = M, \qquad \pi_1 = 0, \qquad \pi_2 = M/2M , \]
de sorte que la fibre de $\mu$ est, en degrés $\leq 2$, le spectre
$\Sigma^2 H(M/2M)$. En appliquant
$R\underline{\mathrm{Hom}}_{H\mathbb{Z}}(-, HN)$ et en tronquant, le troisième
terme du triangle devient
\[ \tau_{\leq 2} \, R\underline{\mathrm{Hom}}(M/2M, N)[-2] \;=\;
   \underline{\mathrm{Hom}}(M/2M, N)[-2] \;=\;
   \underline{\mathrm{Hom}}(M, {}_2 N)[-2] . \]
C'est $(T)$.\footnote{On notera au passage pourquoi $(S)$ n'est pas
manifestement scindée. Le morphisme d'unité $HM \to H\mathbb{Z} \wedge HM$
scinde $\mu$, mais il n'est pas $H\mathbb{Z}$-linéaire : la structure de
module de la source est celle de $HM$, celle du but celle du facteur de
gauche, et elles ne coïncident pas. Sur le topos ponctuel le scindage existe
tout de même, tout complexe de groupes abéliens étant somme de ses
cohomologies ; c'est ce qui rend la question intéressante sur un topos
général, où les faisceaux abéliens ne forment pas une catégorie héréditaire.}

\pagerange{4}{4}
\section*{La suite exacte $(S)$ (page 4)}

Écrivons la suite exacte longue de cohomologie de $(T)$. Le troisième terme
n'a de cohomologie qu'en degré $2$, où elle vaut
$\underline{\mathrm{Hom}}(M, {}_2N)$. Il vient donc, pour $i \leq 1$,
\[ \underline{H}^0(E') = \underline{\mathrm{Hom}}(M,N), \qquad
   \underline{H}^1(E') = \underline{\mathrm{Ext}}^1(M,N), \]
et en degré $2$ la suite exacte
\[ (S) \qquad 0 \longrightarrow \underline{\mathrm{Ext}}^2(M,N)
   \longrightarrow \underbrace{\underline{H}^2\bigl(E'(M,N)\bigr)}_{P(M,N)}
   \xrightarrow{\ \sigma\ } \underline{\mathrm{Hom}}(M, {}_2 N)
   \longrightarrow 0 . \]
L'exactitude à droite est gratuite, et c'est la troncature qui la donne :
$\underline{H}^3(E) = 0$ parce que $E$ est tronqué en degré
$2$.\footnote{Telle que la transcription la lit, la page 4 dit que
$E'(M,N)$ « est un complexe dont les invariants $\underline{H}^i$ sont ceux de
$E(M,N)$ en degré $i+2$ ». Pris à la lettre cela affirmerait
$\underline{H}^i(E') = \underline{H}^{i+2}(E)$, ce qui est faux : le membre de
droite est nul dès que $i > 0$. L'énoncé exact est celui donné ci-dessus, et
c'est manifestement celui que la suite de la page emploie, puisqu'elle enchaîne
aussitôt sur $(S)$.}

La lecture est celle-ci, et c'est tout le contenu de la lettre : classer un
champ de Picard épinglé par $M$ et $N$, c'est classer un champ strict --- une
classe dans $\underline{\mathrm{Ext}}^2(M,N)$ --- et choisir un signe. Rien
d'autre ne se cache entre les deux.

Grothendieck ajoute que $(S)$ « se construit en tout cas canoniquement à la
main », indépendamment de $(T)$ : le terme médian est le faisceau des champs
de Picard épinglés à équivalence près, la flèche $\sigma$ est l'invariant de
symétrie construit plus haut, l'injectivité à gauche est la caractérisation
des champs stricts par $\sigma = 0$, et la surjectivité à droite est l'énoncé
de la section suivante. Cette construction directe ne suppose rien de la
théorie des $2$-champs ; c'est ce qui rend la lettre solide là même où elle se
dit heuristique.\footnote{La page 5 porte, barré, le mot « catégories »,
corrigé en « champs » : la distinction entre les deux --- une catégorie
monoïdale, et un champ dont les sections locales forment de telles catégories
--- est précisément ce qui donne à $(S)$ son terme $\underline{\mathrm{Ext}}^2$
plutôt qu'un $\mathrm{Ext}^2$ global.}

\pagerange{5}{5}
\section*{Tout signe se réalise, et l'obstruction (page 5)}

L'exactitude de $(S)$ comme suite de faisceaux ne dit pas que ses sections se
relèvent : une suite exacte courte de faisceaux abéliens n'est pas exacte sur
les sections. Grothendieck énonce le relèvement séparément, et sous sa forme
la plus forte :

\begin{quote}
\textbf{Énoncé (page 5).} Tout homomorphisme $M \to {}_2N$ provient d'un champ
de Picard épinglé par $M$ et $N$. Par fonctorialité en $X$, il en va de même
au-dessus de tout objet $U$ de $X$ : toute section de
$\underline{\mathrm{Hom}}(M, {}_2N)$ sur $U$ se relève en une section de
$P(M,N)$ sur $U$, et même en un élément de l'hypercohomologie
$\mathbb{H}^2\bigl(U, E'(M,N)\bigr)$.
\end{quote}

La dernière précision est la bonne, et c'est celle qu'il appelle « il y a
mieux en fait » : un élément de $\mathbb{H}^2(U, E')$ est un champ de Picard
véritable sur $U$, alors qu'une section de $P(M,N)$ n'est qu'une donnée locale
recollée à équivalence près.

L'obstruction se lit sur la suite exacte longue d'hypercohomologie de $(T)$ :
\[ \mathbb{H}^2\bigl(U, E'(M,N)\bigr) \longrightarrow
   \mathrm{Hom}_U(M, {}_2N) \xrightarrow{\ \partial\ }
   \mathbb{H}^3\bigl(U, E(M,N)\bigr), \]
et $\mathbb{H}^3(U, E(M,N))$ s'injecte dans $\mathrm{Ext}^3(U; M,
N)$.\footnote{La page dit « a priori l'obstruction est dans
$\mathrm{Ext}^3(X;M,N)$ », en notant $\mathrm{Ext}^n(X;-,-)$ le $\mathrm{Ext}$
global, par opposition au $\underline{\mathrm{Ext}}^n$ faisceautique. C'est
correct mais un peu large : le groupe où vit réellement l'obstruction est
$\mathbb{H}^3(X, \tau_{\leq 2} R\underline{\mathrm{Hom}}(M,N))$, qui s'injecte
dans $\mathrm{Ext}^3(X;M,N)$, le tronqué supérieur $\tau_{\geq 3}$ n'ayant
pas d'hypercohomologie en degré $2$. La précision est la nôtre.} L'énoncé revient donc à
la nullité de $\partial$, que Grothendieck attribue à « un argument
universel » sans le donner.\footnote{Le mot « universel » est entre guillemets
sur la page, et rien de plus n'y figure. La reconstruction naturelle, qui est
la nôtre, consiste à ramener l'énoncé au cas où $M$ est libre. Là le champ
cherché se fabrique à la main : on prend le champ scindé $M \times BN$, dont la
contrainte d'associativité est triviale, et on lui donne pour contrainte de
commutativité une forme biadditive antisymétrique $b$ définie sur une base
$(e_i)$ de $M$ par $b(e_i,e_i) = q(e_i)$ et $b(e_i,e_j) = 0$ pour $i \neq j$.
Sa diagonale vaut bien $q$ --- les valeurs étant d'ordre $2$, $m_i^2$ et $m_i$
y coïncident --- donc $\sigma$ vaut $q$ ; le cas général s'en déduit par une
résolution. Nous ne prétendons
pas que ce soit l'argument qu'il avait en tête.}

Il en tire enfin la formulation qu'il retient, et qui est notablement plus
prudente que celle qu'il avait commencé à écrire : la suite $(S)$ est « bien
proche d'être scindée » --- l'extension admet une section \emph{ensembliste},
au sens où toute section du troisième terme se relève --- sans qu'il soit
affirmé qu'elle se scinde comme suite de faisceaux abéliens.\footnote{Le mot
« splittée » est écrit puis barré sur la page. La prudence est justifiée : le
relèvement de toutes les sections au-dessus de tout objet est strictement plus
faible qu'un scindage additif, et l'on a vu plus haut pourquoi le scindage
évident échoue sur un topos général.}

\pagerange{6}{6}
\section*{L'exemple : $K^0$ et $K^1$ (page 6)}

Soit $A$ un anneau sur $X$, et soit considéré le champ additif des
$A$-modules projectifs, muni de la somme directe. Prenons pour $M$ et $N$ les
faisceaux $K^0$ et $K^1$ qui lui sont associés. Alors « le raisonnement de Mme Sinh »
fournit un champ de Picard épinglé par $M$ et $N$, d'où une section canonique
du terme médian $P(M,N)$ de $(S)$.\footnote{C'est le contenu du dossier 135 :
l'enveloppe de Picard --- aujourd'hui : la complétion en groupe --- du champ
additif des modules projectifs a pour invariants $K^0$ et $K^1$. La lettre
écrit « $L \otimes L$ » pour la loi monoïdale générique ; ici cette loi est la
somme directe, et c'est $L \oplus L$ qu'il faut lire dans ce paragraphe.}

L'invariant de symétrie de ce champ se calcule, et il vaut la peine de
l'écrire parce qu'il est complètement explicite. Pour $L$ un module projectif,
$c_{L,L}$ est l'automorphisme d'échange de $L \oplus L$, c'est-à-dire la
matrice de permutation $\left(\begin{smallmatrix} 0 & 1 \\ 1 &
0\end{smallmatrix}\right)$ par blocs. Sa classe dans $K^1$ est celle de
$-1$ : pour $L$ de rang $r$, l'échange est un produit de $r$ transpositions,
donc de déterminant $(-1)^r$. Autrement dit
\[ \sigma \,:\, K^0 \longrightarrow {}_2 K^1, \qquad
   x \longmapsto \{-1\} \cdot x . \]
Dans le langage des spectres, c'est la multiplication par l'élément $\eta$ de
$\pi_1(\mathbb{S})$, dont l'image dans $K^1$ est précisément $\{-1\}$. La
K-théorie ne fournit donc pas seulement une section de $(S)$ : elle fournit
\emph{la} section universelle, celle que porte la sphère.\footnote{Cette
identification est la nôtre. La lettre s'arrête à l'existence de la section
canonique et n'en calcule pas $\sigma$. Le calcul par la matrice d'échange
est élémentaire ; l'énoncé « $\eta$ s'envoie sur $\{-1\}$ » est un fait
standard de K-théorie algébrique, postérieur en tout cas à la mise en forme
de la lettre.}

Grothendieck note enfin que tout ce qui précède a « les fonctorialités
évidentes en $M$, $N$, $X$ ».

\pagerange{6}{6}
\section*{La question, et ce qu'on en sait aujourd'hui (page 6)}

La lettre s'achève sur ce pour quoi elle est écrite :

\begin{quote}
Le triangle exact $(T)$ et la suite exacte $(S)$ sont-ils connus par les
compétents (Quillen, Breen, Illusie…) ? Connaissent-ils des variantes
« supérieures » ? Un principe « géométrique » pour les obtenir pourrait être
via des $n$-champs de Picard non nécessairement stricts.
\end{quote}

Le vocabulaire d'aujourd'hui permet de dire où se rangent ces deux
questions, sans rien dire de qui savait quoi et quand.\footnote{Cette lecture
ne répond pas à la question posée, qui est une question de fait sur l'état des
connaissances vers 1975 et que les seuls documents du dossier ne tranchent
pas. Elle se borne à nommer les objets. On ne trouvera pas non plus ici
d'affirmation sur l'antériorité : la lettre demande si la chose est connue,
ce qui est le contraire d'une revendication.}

Un champ de Picard sur $X$ est un faisceau de spectres connectifs tronqué en
degré $1$ ; ses invariants sont $\pi_0 = M$, $\pi_1 = N$, et sa classification
est celle des tours de Postnikov à deux étages, le premier invariant vivant
dans le groupe des opérations stables de degré $2$ de $M$ vers $N$. Ce groupe
est exactement le $\underline{H}^2(E')$ de la lettre. Les champs stricts sont
ceux qui sont des modules sur $H\mathbb{Z}$, et leur classification par
$\underline{\mathrm{Ext}}^2$ est le théorème de Deligne. Le triangle $(T)$ est
alors la comparaison entre les deux linéarités, $\mathbb{Z}$ et
$\mathbb{S}$ ; l'écart, en degré $2$, est le dual du carré de Steenrod
$Sq^{2}$, et c'est la multiplication par $\eta$.

Quant aux « variantes supérieures », la réponse de principe est oui, et elle
est celle que Grothendieck devine en écrivant « via des $n$-champs de Picard
non nécessairement stricts » : un $n$-champ de Picard est un faisceau de
spectres connectifs tronqué en degré $n$, les stricts sont les
$H\mathbb{Z}$-modules, et la comparaison générale entre
$R\underline{\mathrm{Hom}}_{\mathbb{S}}$ et
$R\underline{\mathrm{Hom}}_{H\mathbb{Z}}$ est gouvernée par l'homologie
entière de $H\mathbb{Z}$ tout entière. Le triangle $(T)$ est le premier étage
de cette tour --- celui que produit $H_2(H\mathbb{Z};\mathbb{Z}) =
\mathbb{Z}/2$ --- et il y en a un par degré au-dessus. C'est la tour de
Postnikov de la sphère qui organise l'ensemble, et le « principe géométrique »
demandé est là.\footnote{Les noms modernes à chercher sont : catégorie ou
champ de Picard, type d'homotopie stable de dimension $1$, invariant de
Postnikov, opérations cohomologiques stables, $\eta$, $Sq^2$, et pour la
version faisceautique les travaux de Larry Breen --- l'un des trois
« compétents » nommés dans la question --- sur les gerbes et les $2$-gerbes.
Aucune de ces références ne figure sur la page, et rien n'indique que
Grothendieck ait connu sous ce nom l'une quelconque d'entre elles ; la mise en
correspondance est la nôtre.}

\end{document}
